% =====================================================================
\section{The sleep switch solved on paper}\label{sec:anal}
The whole of section~\ref{sec:bifur} was obtained by computation. The question then arose whether the same could be had on paper, and it can: with the fast pair of voltages eliminated adiabatically, the switch is a relay with hysteresis, and its return map admits a closed form up to two transcendental equations. The reduction settles, among other things, the question that the scan could not settle --- whether there are invariant tori in the window.

\textbf{The folds in closed form.} The steady state of the fast pair admits a \emph{parametric} solution: taking $V_m$ as the parameter,
\[Q(V_v)=\frac{A_m-V_m}{|\nu_{mv}|},\qquad D_v(V_m)=V_v-\nu_{vm}Q(V_m),\]
with $Q^{-1}(q)=\theta+\sigma_Q\ln\bigl(q/(Q_{\max}-q)\bigr)$, and the folds are the extrema of $D_v(V_m)$. Evaluating: $1.4503$ and $2.4635$\,mV against the $1.450$ and $2.463$ quoted by the computation, to four significant figures \emph{without integration}.

\textbf{The return map.} In the relay limit the somnogens obey linear equations with closed-form solutions, and switching occurs at two \emph{moving thresholds} $g_\pm(t)$. If $t_n$ is a sleep onset, then $t_{n+1}=f(t_n)$ is given implicitly by
\[g_+(t_n)\,e^{-(t_w-t_n)/\tau_s}=g_-(t_w),\qquad 1-\bigl(1-g_-(t_w)\bigr)e^{-(t_{n+1}-t_w)/\tau_w}=g_+(t_{n+1}),\]
both of the form ``exponential $=$ sinusoid''. The free-running period follows in closed form: $T_s=1.337$\,h, $T_w=6.683$\,h, $T_0=8.020$\,h against $8.878$\,h from direct integration. The 10\,\% discrepancy is not an error of derivation: switching does \emph{not} occur at the fold. Sleep onset is observed at $D_v=2.5026$ against the fold at $2.4635$, because the relaxation time diverges at a saddle-node and a slow passage completes after the bifurcation point. The overshoot was measured over two decades in the fast time scale and vanishes with exponent $0.64$ and $0.59$, against the $2/3$ expected for slow passage through a saddle-node. The closed form is therefore the exact relay limit, and the model stands 10\,\% away from it --- which must be stated rather than concealed.

\textbf{The tongue edge without a scan.} The $1{:}1$ orbit is a fixed point of $f$ shifted by one day, and the edge is its saddle-node: $f(t^*)=t^*+24$ together with $f'(t^*)=1$ gives three equations in three unknowns. The result is $4.7651$\,mV; the independent scan of the rotation number gave $4.77$\,mV.

\textbf{Arnold tongues do not apply at all.} The somnogens rise towards unity from below, so a sleep onset is possible only where $g_+(t)<1$. The fraction of the day admitting one falls from 0.789 at zero coupling to 0.166 at $4.80$\,mV and to 0.016 at $5.18$, where it vanishes --- and that is how sleep is lost. At $4.80$\,mV the domain is an arc of length $3.97$\,h and its image an arc of length $0.993$ inside it. A degree-one circle map has an image covering the whole circle; here the image occupies 4\,\% of it. The map is piecewise smooth with discontinuities: at $4.88$\,mV the jump is $22.47$\,h against a median step of $10^{-4}$\,h --- almost exactly one day, that is, a skipped night. \emph{The locking staircase is generated by branch switching, not by mode locking of a smooth circle map.}

\textbf{No invariant tori, and this is an inequality rather than an observation.} The right quantity is the Lyapunov exponent \emph{along the orbit}. Iterating the map and discarding the transient, over the whole window $\lambda$ lies between $-3.54$ and $-1.62$, with rotation numbers $1$, $3/2$, $2$, $3$ reproducing the staircase. Two conclusions follow: chaos is excluded, since it requires $\lambda>0$; and a smooth invariant torus is excluded, since a quasiperiodic attractor requires $\lambda=0$ exactly, whereas the distance to zero here is at least $1.6$. The attractor at every parameter examined is a strongly attracting periodic orbit. This is the third independent confirmation of the reduction, after the folds and the tongue edge.

\begin{note}
Reaching the right quantity required discarding two wrong ones. The first was the locked fraction on a grid: a locked fraction below unity was taken to imply that tori occur on a parameter set of positive measure. The inference is invalid --- the complement of the locked set is \emph{dense}, so a finite grid finds unlocked nodes whether the complement has positive measure or measure zero, and misses them in both cases with a probability depending only on the resolution. The criterion has no power in either direction, and the honest content of ``240 of 240 nodes locked'' is a lower bound on the measure of the locked set, $0.988$ at 95\,\% confidence, and nothing more. The second wrong quantity was the maximum of $|f'|$ over the domain: at tangential crossings the derivative has a square-root singularity, and a grid samples points the attractor never visits. The general lesson: on a Poincar\'e section, a quantity averaged \emph{along the orbit} and a quantity taken over the \emph{whole domain} answer different questions, and the latter need not concern the attractor at all.
\end{note}

\begin{ineq}
Measure in \emph{phase} space is a different question and does not arise here. All of the above concerns parameter space. The measure of the attractor itself --- a thick attractor of positive Lebesgue measure, whose dimension equals that of the system --- belongs to a harder circle of problems; Ilyashenko constructed examples for step skew products. The present system sits at the opposite pole, and this is measured rather than assumed: the fast time scales (10\,s) against the slow ones (hours) give a strongly negative divergence, and the two fast dimensions collapse at once --- the control established one-dimensionality of the return map to $4.5\cdot10^{-13}$, eight trajectories from different initial somnogen values falling onto a single curve to machine precision. The one-dimensional map then contracts onto a periodic orbit with $\lambda\le-1.6$, so the attractor on the section is a finite set of points and, in the extended phase space, a closed curve. This is structural: approaching thickness would require the divergence to cease being negative somewhere, that is, the fast time scales to become comparable to the slow ones --- but the entire reduction to a relay, and with it the folds and the circle map, rests on their separation. What formally remains open is an exceptional parameter carrying a Denjoy-type Cantor minimal attractor; the argument against is again structural rather than statistical, since such an attractor also requires the exponent at zero.
\end{ineq}

\textbf{What this buys in computation.} Constructing the map requires two scalar roots per cycle instead of 1600 integration steps per day, and the tongue edge requires solving three equations instead of a bisection. The scan took 1686\,s; the analytic route yields the same numbers in seconds, up to the stated 10\,\% offset. It applies in particular to the part the scan returned as untestable: the exponent along the orbit is obtained in twelve iterations of the map and answers its question directly, whereas a grid of 240 nodes could not answer it in principle, since it searched for aperiodic nodes where the attractor is periodic at every parameter.
